\begin{abstract}
Agricultural irrigation accounts for approximately 70\% of global freshwater withdrawals, with an estimated 40\% to 60\% squandered due to inefficient open-calendar flood irrigation schedules and decoupled thresholding heuristics. While cloud-centric deep learning frameworks have demonstrated high diagnostic accuracy, their reliance on power-intensive servers, high-bandwidth telemetry, and multi-million-parameter neural networks renders them fundamentally unsuitable for remote agricultural regions characterized by cellular dead zones and stringent power limitations. Conversely, existing edge-based microcontrollers frequently depend on costly thermal radiometric imagers, unengineered sensor statistics, and unconstrained decision trees that induce class collapse under severe moisture deficit states. In this paper, we propose \textbf{HydroEdge-Net}, an explainable, biophysically informed, ultra-lightweight TinyML edge computing framework designed for autonomous crop water stress diagnosis and closed-loop irrigation scheduling on low-cost microcontrollers. HydroEdge-Net couples microclimatic and edaphic telemetry via Tetens-derived Vapor Pressure Deficit (VPD), Soil Moisture Deficit (SMD) relative to field capacity ($\theta_{\text{FC}} = 35.0\%$), and a novel Thermal-Moisture Stress Coupling Index ($\Psi_{\text{TMS}}$), augmented with categorical soil texture and phenological stage encodings. A cost-sensitively reweighted, depth-constrained gradient boosted decision tree ($d \le 5$, 45 estimators) is deployed via a bare-metal static C header without heap allocation or external runtime dependencies. Evaluated on a multi-crop benchmark of 16,411 instances across five crop varieties and seven soil textures, HydroEdge-Net achieves 95.71\% overall accuracy (99.18\% in tuned configuration) and an exceptional 98.96\% water deficit recall. When compiled for an Espressif ESP32-S3 microcontroller (240~MHz), the model occupies merely 16.8~KB of Flash and 11.4~KB of SRAM with an active execution latency of 1.42~ms, demonstrating a 40.0\% Flash reduction and 64.38\% SRAM reduction relative to state-of-the-art edge baselines (AquaSense). Furthermore, local TreeSHAP explainability confirms rigorous alignment with physiological transpiration dynamics. In a 30-day (720-hour) dual-layer FAO-56 Penman-Monteith hydrological digital twin simulation, HydroEdge-Net achieves a 63.96\% water conservation gain over calendar schedules while completely preventing transpirational water stress and root-zone hypoxia, whereas static thresholding suffered recurrent crop moisture deficits.
\end{abstract}

\begin{IEEEkeywords}
Precision agriculture, TinyML, edge artificial intelligence, crop water stress, explainable AI (XAI), Soil-Plant-Atmosphere Continuum, autonomous irrigation, FAO-56 Penman-Monteith, digital twin.
\end{IEEEkeywords}
